\documentclass[10pt,a4paper]{amsart}
\usepackage[margin=19mm]{geometry}
\usepackage{amsmath,amssymb,mathtools}
\usepackage[scr=rsfs,cal=euler]{mathalfa}
\usepackage{xfrac}
\usepackage[unicode,hidelinks,bookmarks=false]{hyperref}
\newcommand{\ie}{i.e.\ }
\newcommand{\cf}{cf.\ }
\newcommand{\resp}{respectively\ }
\newcommand{\Subsection}{\subsection}
\providecommand{\nobreakdash}{\nobreak\hbox{-}}
\setlength{\emergencystretch}{2em}
\multlinegap0pt
\usepackage{bm}
\usepackage{bbm}
\usepackage{dsfont}
\usepackage{xspace}
\usepackage{cancel}
\usepackage{mathtools}
\usepackage{amsthm}
\makeatletter
\@ifundefined{theorem}{\newtheorem{theorem}{Theorem}}{}
\@ifundefined{lemma}{\newtheorem{lemma}[theorem]{Lemma}}{}
\makeatother
\newcommand{\R}{\mathbb{R}}
\newcommand{\dd}{\,\mathrm{d}}
\title[Bifurcation analysis of Stokes waves with piecewise smooth v]{Bifurcation analysis of Stokes waves with piecewise smooth vorticity in deep water}
\author{Changfeng Gui, Jun Wang, Wen Yang, Yong Zhang}
\date{Source: arXiv:2511.03973v3 (CC BY 4.0)}
\begin{document}
\maketitle

\noindent\emph{Source and licence.} Bifurcation analysis of Stokes waves with piecewise smooth vorticity in deep water. Version arXiv:2511.03973v3 (CC BY 4.0), distributed under the Creative Commons Attribution 4.0 International licence, as recorded in the arXiv licence metadata for this exact version and shown on the abstract page https://arxiv.org/abs/2511.03973v3. Full rights notice: licenses/arxiv-2511.03973-CC-BY-4.0.txt.\par\smallskip

\noindent\textit{Editorial scope.} Counted unit: the Lemma whose source label is \texttt{lem4.3} in the article (source0.tex lines 1347--1433, source label \texttt{lem4.3}), delivered together with the article's own complete original proof of it (source0.tex lines 1435--1861, one \texttt{proof} environment of 427 lines). No mathematics is written, repaired or re-derived: statement and proof are the article's own text line for line. Required context retained uncounted: eq2.bifcondition (lines 500--507); eq:quadratic-form (lines 1319--1329); eq:principal-Rayleigh-value (lines 1338--1345). Every definition, display and label of those lines is retained unchanged. Omitted from this package and not redistributed: 1--499, 508--1318, 1330--1337, 1346--1346, 1434--1434, 1862--5453. Nothing inside a retained excerpt is omitted or altered: every retained body is delivered exactly as the article's lines are written, so no omission receipt of any kind accompanies this package. Citations: the retained excerpts carry no citation command at all, so no bibliography entry of the article is transcribed and no reference is redistributed. Reference adaptation: none was needed and none was made; every reference inside the delivered unit resolves inside it, so no unresolved reference or citation is produced. Printed slips kept literal: no printed word, symbol, hypothesis or number of the counted statement or of its proof was altered. Layout and class: the article's own class is \texttt{amsart} with packages including mathrsfs, amsmath, amssymb, color, amsthm, epsfig, graphicx, titletoc, fullpage, palatino, mathpazo; the wrapper is the lane's pinned \texttt{amsart} editorial wrapper at 10pt on A4, which restates the theorem-like environments the retained text needs and adds the article's own macro definitions that the retained text needs (\texttt{\textbackslash R}, \texttt{\textbackslash dd}), with extra wrapper packages (bm, bbm, dsfont, xspace, cancel, mathtools). The delivered numbering is the unit's own. Assets: no retained span contains a figure environment, a graphics-inclusion command, a PDF-page inclusion, a table or a TikZ picture; no image file is copied into this package, the package declares no asset dependency and no image file is redistributed with it. Licence: version arXiv:2511.03973v3 of this article is distributed under the Creative Commons Attribution 4.0 International licence, as recorded in the arXiv licence metadata for this exact version and shown on the abstract page https://arxiv.org/abs/2511.03973v3. A full rights notice is delivered at licenses/arxiv-2511.03973-CC-BY-4.0.txt. Characters: every retained excerpt is pure ASCII, so the delivered engine, XeLaTeX with its default Latin Modern text font, reports no missing character.
\begin{equation}\label{eq2.bifcondition}
\int_{-\infty}^{0}
\left\{
2\bigl(2\Gamma(p)-2\Gamma_{\inf}\bigr)^{3/2}
+
\bigl(2\Gamma(p)-2\Gamma_{\inf}\bigr)^{1/2}
\right\}e^{2p}\,\dd p<g.
\end{equation}

\begin{equation}
\label{eq:quadratic-form}
\mathfrak q_{\varepsilon,\lambda}[\Phi]
:=
-g\Phi(0)^2
+
\int_{-\infty}^{0}
\left(
a^3(p;\lambda)\Phi_p^2+\varepsilon\Phi^2
\right)\dd p
\end{equation}

\begin{equation}
\label{eq:principal-Rayleigh-value}
\mu^\varepsilon(\lambda)
:=
\inf_{\substack{\Phi\in H^1(-\infty,0)\\ \Phi\not\equiv 0}}
\frac{\mathfrak q_{\varepsilon,\lambda}[\Phi]}
{\|\Phi\|_{a,\lambda}^2}.
\end{equation}

\begin{lemma}[Kernel of the linearized approximate problem]
\label{lem4.3}
Assume that
\[
\gamma\in
C^{1,\alpha}\bigl([0,-p_0)\bigr)
\cap
C^{1,\alpha}\bigl([-p_0,\infty)\bigr),
\qquad
\alpha\in(0,1),
\]
that
\[
\gamma(s)=O(s^{-2-r})
\qquad\text{as }s\to\infty
\]
for some $r>0$, and that the local-bifurcation condition \eqref{eq2.bifcondition} holds.

Then there exists $\varepsilon_0>0$ such that, for every
$0<\varepsilon<\varepsilon_0$, there is a unique
\[
\lambda_*^\varepsilon\in(-2\Gamma_{\inf},\infty)
\]
for which
\begin{equation}
\label{eq:crossing}
\mu^\varepsilon(\lambda_*^\varepsilon)=-1.
\end{equation}
At $\lambda=\lambda_*^\varepsilon$, the kernel of the linearized operator is
one-dimensional:
\begin{equation}
\label{eq:kernel-formula}
\ker
\partial_{(w,\underline w)}
F^\varepsilon(\lambda_*^\varepsilon,0,0)
=
\operatorname{span}
\left\{
\left(
\phi^\varepsilon(p)\cos q,\,
\varphi^\varepsilon(p)\cos q
\right)
\right\}.
\end{equation}
Here
\[
\Psi^\varepsilon(p)
=
\begin{cases}
\phi^\varepsilon(p),&p\in[p_0,0],\\
\varphi^\varepsilon(p),&p\in(-\infty,p_0],
\end{cases}
\]
is a nontrivial eigenfunction satisfying
\begin{equation}
\label{eq:one-dimensional-eigenproblem}
\left\{
\begin{aligned}
-\bigl(a^3(p;\lambda_*^\varepsilon)\Psi_p^\varepsilon\bigr)_p
+\varepsilon\Psi^\varepsilon
&=
-a(p;\lambda_*^\varepsilon)\Psi^\varepsilon
&&\text{in }(-\infty,p_0)\cup(p_0,0),\\
(\lambda_*^\varepsilon)^{3/2}\Psi_p^\varepsilon(0)
&=g\Psi^\varepsilon(0),\\
[\Psi^\varepsilon]_{p=p_0}
=
[\Psi_p^\varepsilon]_{p=p_0}
&=0,\\
\Psi^\varepsilon(p),\ \Psi_p^\varepsilon(p)
&\longrightarrow0
&&\text{as }p\to-\infty.
\end{aligned}
\right.
\end{equation}
Moreover,
\begin{equation}
\label{eq:root-convergence}
\lambda_*^\varepsilon\longrightarrow\lambda_*^0
\qquad\text{as }\varepsilon\downarrow0,
\end{equation}
where $\lambda_*^0$ is the unique solution of
\[
\mu^0(\lambda)=-1
\]
in $(-2\Gamma_{\inf},\infty)$.
\end{lemma}

\begin{proof}
We divide the proof into five steps.

\medskip
\noindent
\textbf{Step 1: Variational realization of every negative Rayleigh value.}

Fix $\lambda>-2\Gamma_{\inf}$ and $\varepsilon\geq0$.  Then
\[
a(p;\lambda)\geq
a_{\min}(\lambda)
:=
\bigl(\lambda+2\Gamma_{\inf}\bigr)^{1/2}>0.
\]
Consequently, the denominator in
\eqref{eq:principal-Rayleigh-value} is equivalent to the
$L^2(-\infty,0)$ norm.  The trace inequality on the half-line gives,
for every $\eta>0$,
\[
|\Phi(0)|^2
\leq
\eta\|\Phi_p\|_{L^2(-\infty,0)}^2
+
C_\eta\|\Phi\|_{L^2(-\infty,0)}^2.
\]
It follows that the form
$\mathfrak q_{\varepsilon,\lambda}$ is closed and bounded from below on
$H^1(-\infty,0)$.

Suppose that $\mu^\varepsilon(\lambda)<0$.  Let $\{\Phi_n\}$ be a minimizing
sequence normalized by
\[
\|\Phi_n\|_{a,\lambda}=1.
\]
The trace estimate and the positive lower bound for $a$ show that
$\{\Phi_n\}$ is bounded in $H^1(-\infty,0)$.  After passing to a
subsequence,
\[
\Phi_n\rightharpoonup\Phi
\quad\text{weakly in }H^1(-\infty,0),
\]
and
\[
\Phi_n\longrightarrow\Phi
\quad\text{strongly in }L^2_{\mathrm{loc}}(-\infty,0]
\]
and at the trace point $p=0$.

The weak limit cannot vanish identically.  Indeed, if $\Phi\equiv0$,
then $\Phi_n(0)\to0$, whereas all integral terms in
\eqref{eq:quadratic-form} are nonnegative.  This would imply
\[
\liminf_{n\to\infty}
\mathfrak q_{\varepsilon,\lambda}[\Phi_n]\geq0,
\]
contrary to $\mu^\varepsilon(\lambda)<0$.

By weak lower semicontinuity,
\[
\mathfrak q_{\varepsilon,\lambda}[\Phi]
\leq
\liminf_{n\to\infty}
\mathfrak q_{\varepsilon,\lambda}[\Phi_n]
=
\mu^\varepsilon(\lambda),
\]
while
\[
0<\|\Phi\|_{a,\lambda}^2\leq1.
\]
Since the right-hand side is negative,
\[
\frac{\mathfrak q_{\varepsilon,\lambda}[\Phi]}
{\|\Phi\|_{a,\lambda}^2}
\leq
\mathfrak q_{\varepsilon,\lambda}[\Phi]
\leq
\mu^\varepsilon(\lambda).
\]
The opposite inequality follows from the definition of
$\mu^\varepsilon(\lambda)$.  Hence the infimum is attained.

The Euler--Lagrange equation is
\begin{equation}
\label{eq:euler-lagrange-general}
-\bigl(a^3\Psi_p\bigr)_p+\varepsilon\Psi
=
\mu^\varepsilon(\lambda)a\Psi
\quad\text{in }(-\infty,p_0)\cup(p_0,0),
\end{equation}
with
\[
a^3(0;\lambda)\Psi_p(0)=g\Psi(0).
\]
Since $\Gamma(0)=0$, we have $a^3(0;\lambda)=\lambda^{3/2}$.
The weak formulation also yields continuity of $\Psi$ and of
$a^3\Psi_p$ at $p=p_0$.  As $a$ is continuous and strictly positive,
this is equivalent to
\[
[\Psi]_{p=p_0}=[\Psi_p]_{p=p_0}=0.
\]
Piecewise elliptic regularity gives
\[
\Psi\in
C^{3,\alpha}\bigl((-\infty,p_0]\bigr)
\cap
C^{3,\alpha}\bigl([p_0,0]\bigr).
\]
Moreover, since the coefficients converge at $p=-\infty$ and the
spectral value is negative, the one-dimensional equation has an
exponential dichotomy.  The $H^1$ condition removes the growing
solution; consequently, for some $C,\rho>0$,
\[
 |\partial_p^j\Psi(p)|\leq Ce^{\rho p},
 \qquad p\leq p_0-1,\qquad 0\leq j\leq3.
\]
In particular the eigenfunction belongs to the function space used in
\eqref{eq:kernel-formula}.

\medskip
\noindent
\textbf{Step 2: The nonpositive spectral subspace has dimension at most one.}

For every nonzero $\Phi\in H^1(-\infty,0)$ satisfying $\Phi(0)=0$,
\[
\mathfrak q_{\varepsilon,\lambda}[\Phi]
=
\int_{-\infty}^{0}
\left(
a^3\Phi_p^2+\varepsilon\Phi^2
\right)\dd p>0
\]
for every $\varepsilon\geq0$.  Therefore, the maximal dimension of a subspace on
which $\mathfrak q_{\varepsilon,\lambda}$ is nonpositive is at most one.
Indeed, if a two-dimensional subspace $E$ were nonpositive, then the
linear trace map
\[
E\ni\Phi\longmapsto\Phi(0)\in\R
\]
would have a nontrivial kernel.  A nonzero element of that kernel
would have strictly positive quadratic form, a contradiction.

It follows in particular that, whenever
$\mu^\varepsilon(\lambda)<0$, this negative eigenvalue is simple and is the
only nonpositive eigenvalue of the generalized Sturm--Liouville
problem \eqref{eq:euler-lagrange-general}.

\medskip
\noindent
\textbf{Step 3: Existence and uniqueness of the crossing
$\mu^\varepsilon(\lambda)=-1$.}

Set
\[
\lambda_0:=-2\Gamma_{\inf},
\qquad
a_0(p):=
\bigl(2\Gamma(p)-2\Gamma_{\inf}\bigr)^{1/2}.
\]
Condition \eqref{eq2.bifcondition} implies, in particular, that
\[
\delta_0
:=
g-
\int_{-\infty}^{0}
\bigl(a_0^3(p)+a_0(p)\bigr)e^{2p}\,\dd p
>0.
\]
Choose
\[
0<\varepsilon_0<\delta_0.
\]
For $\lambda>\lambda_0$, evaluate the Rayleigh quotient at
$\Phi(p)=e^p$.  Since
\[
\int_{-\infty}^{0}e^{2p}\,\dd p=\frac12,
\]
we have
\begin{align}
&\mathfrak q_{\varepsilon,\lambda}[e^p]
+\|e^p\|_{a,\lambda}^2
\nonumber\\
&\qquad=
-g+
\int_{-\infty}^{0}
\left(
a^3(p;\lambda)+a(p;\lambda)
\right)e^{2p}\,\dd p
+\frac{\varepsilon}{2}.
\label{eq:left-trial-expression}
\end{align}
As $\lambda\downarrow\lambda_0$, the right-hand side of
\eqref{eq:left-trial-expression} converges to
\[
-\delta_0+\frac{\varepsilon}{2}.
\]
After decreasing $\varepsilon_0$ if necessary, there exists
$\lambda_->\lambda_0$, independent of
$0<\varepsilon<\varepsilon_0$, such that
\[
\mathfrak q_{\varepsilon,\lambda_-}[e^p]
+\|e^p\|_{a,\lambda_-}^2<0.
\]
Hence
\begin{equation}
\label{eq:left-sign}
\mu^\varepsilon(\lambda_-)<-1.
\end{equation}

Next set
\[
\lambda_+:=g-2\Gamma_{\inf}.
\]
For every $\Phi\in H^1(-\infty,0)$, we have
\[
a(p;\lambda_+)\geq\sqrt g.
\]
Moreover,
\[
\Phi(0)^2
=
2\int_{-\infty}^{0}\Phi\Phi_p\,\dd p
\leq
g^{-1/2}\int_{-\infty}^{0}\Phi^2\,\dd p
+
g^{1/2}\int_{-\infty}^{0}\Phi_p^2\,\dd p.
\]
Consequently,
\begin{align*}
\mathfrak q_{\varepsilon,\lambda_+}[\Phi]
+\|\Phi\|_{a,\lambda_+}^2
&=
-g\Phi(0)^2
+
\int_{-\infty}^{0}
\left(
a^3\Phi_p^2+(a+\varepsilon)\Phi^2
\right)\dd p\\
&\geq
\varepsilon\int_{-\infty}^{0}\Phi^2\,\dd p>0
\end{align*}
for every nonzero $\Phi$ when $\varepsilon>0$.  Since
$a(\cdot;\lambda_+)$ is bounded above, with
\[
 a_{\max,+}:=\sup_{p\leq0}a(p;\lambda_+)<\infty,
\]
we have
\[
 \int_{-\infty}^0\Phi^2\,\dd p
 \geq a_{\max,+}^{-1}\|\Phi\|_{a,\lambda_+}^2.
\]
Therefore,
\begin{equation}
\label{eq:right-sign}
\mu^\varepsilon(\lambda_+)
\geq-1+\frac{\varepsilon}{a_{\max,+}}>-1.
\end{equation}

On every compact interval contained in
$(-2\Gamma_{\inf},\infty)$, the coefficients of the quadratic forms
depend continuously on $\lambda$, and hence
$\lambda\mapsto\mu^\varepsilon(\lambda)$ is continuous.  The signs
\eqref{eq:left-sign}--\eqref{eq:right-sign} therefore yield at least
one
\[
\lambda_*^\varepsilon\in(\lambda_-,\lambda_+)
\]
such that \eqref{eq:crossing} holds.

We now prove uniqueness.  Whenever $\mu^\varepsilon(\lambda)<0$, Step~2
shows that it is a simple isolated eigenvalue.  We normalize its
eigenfunction by
\[
\int_{-\infty}^{0}a(p;\lambda)\Psi^2(p)\,\dd p=1.
\]
Standard perturbation theory for simple eigenvalues of closed
quadratic forms implies differentiability with respect to $\lambda$.
Since
\[
\partial_\lambda a=\frac1{2a},
\qquad
\partial_\lambda(a^3)=\frac32a,
\]
the Hellmann--Feynman formula gives
\begin{equation}
\label{eq:mu-derivative}
\frac{\dd}{\dd\lambda}\mu^\varepsilon(\lambda)
=
\frac32
\int_{-\infty}^{0}a\Psi_p^2\,\dd p
-
\frac{\mu^\varepsilon(\lambda)}2
\int_{-\infty}^{0}a^{-1}\Psi^2\,\dd p.
\end{equation}
Both terms on the right-hand side are strictly positive whenever
$\mu^\varepsilon(\lambda)<0$.  Thus
\[
\frac{\dd}{\dd\lambda}\mu^\varepsilon(\lambda)>0
\qquad\text{whenever }\mu^\varepsilon(\lambda)<0.
\]
In particular, the equation
$\mu^\varepsilon(\lambda)=-1$ has exactly one solution.  This
conclusion follows from the strict derivative formula
\eqref{eq:mu-derivative}; no unproved global shape of the eigenvalue
curve is being assumed.
\medskip
\noindent
\textbf{Step 4: Identification of the full kernel, including the
zero Fourier mode.}

Let $(u,\underline u)$ belong to the kernel of
$\partial_{(w,\underline w)}F^\varepsilon(\lambda_*^\varepsilon,0,0)$.  Since the
function space consists of even $2\pi$-periodic functions, we must
include the zero Fourier mode and write
\[
u(q,p)=\sum_{k=0}^{\infty}\phi_k(p)\cos(kq),
\qquad
\underline u(q,p)
=
\sum_{k=0}^{\infty}\varphi_k(p)\cos(kq).
\]
For each $k\geq0$, the piecewise function
\[
\Psi_k(p)
=
\begin{cases}
\phi_k(p),&p\in[p_0,0],\\
\varphi_k(p),&p\in(-\infty,p_0],
\end{cases}
\]
satisfies
\begin{equation}
\label{eq:k-mode}
-\bigl(a^3\Psi_{k,p}\bigr)_p+\varepsilon\Psi_k
=
-k^2a\Psi_k,
\qquad
\lambda^{3/2}\Psi_{k,p}(0)=g\Psi_k(0),
\end{equation}
together with the transmission and decay conditions.

At $\lambda=\lambda_*^\varepsilon$, the value $-1$ is the principal
eigenvalue.  By Step~2 it is simple and is the only nonpositive
eigenvalue.  The mode $k=0$ would require the eigenvalue $0$, while
every mode $k\geq2$ would require the eigenvalue
$-k^2<-1$.  Both possibilities are impossible.  Therefore only
$k=1$ occurs, and its one-dimensional eigenspace yields
\eqref{eq:kernel-formula}.

\medskip
\noindent
\textbf{Step 5: Convergence of the entire family
$\lambda_*^\varepsilon$ as $\varepsilon\downarrow0$.}

Let
\[
I:=[\lambda_-,\lambda_+]
\Subset(-2\Gamma_{\inf},\infty)
\]
and set
\[
a_*:=\inf_{\substack{\lambda\in I\\p\leq0}}
a(p;\lambda)
=
\bigl(\lambda_-+2\Gamma_{\inf}\bigr)^{1/2}>0.
\]
For every nonzero $\Phi\in H^1(-\infty,0)$ and every $\lambda\in I$,
\[
0
\leq
\frac{\mathfrak q_{\varepsilon,\lambda}[\Phi]}
{\|\Phi\|_{a,\lambda}^2}
-
\frac{\mathfrak q_{0,\lambda}[\Phi]}
{\|\Phi\|_{a,\lambda}^2}
=
\varepsilon
\frac{\displaystyle\int_{-\infty}^{0}\Phi^2\,\dd p}
{\displaystyle\int_{-\infty}^{0}a\Phi^2\,\dd p}
\leq
\frac{\varepsilon}{a_*}.
\]
Taking infima gives
\begin{equation}
\label{eq:uniform-mu-convergence}
0
\leq
\mu^\varepsilon(\lambda)-\mu^0(\lambda)
\leq
\frac{\varepsilon}{a_*}
\qquad
\text{for every }\lambda\in I.
\end{equation}
Thus $\mu^\varepsilon\to\mu^0$ uniformly on $I$.

The argument above for $\varepsilon=0$ still shows that every negative
Rayleigh value is attained, simple, and strictly increasing in
$\lambda$.  Moreover,
\[
\mu^0(\lambda_-)<-1,
\qquad
\mu^0(\lambda_+)\geq-1.
\]
Hence there exists a unique
$\lambda_*^0\in I$ satisfying
\[
\mu^0(\lambda_*^0)=-1.
\]

Let $\varepsilon_j\downarrow0$ and suppose that, along a subsequence,
$\lambda_*^{\varepsilon_j}\to\bar\lambda\in I$.  By
\eqref{eq:uniform-mu-convergence},
\[
-1
=
\mu^{\varepsilon_j}(\lambda_*^{\varepsilon_j})
\longrightarrow
\mu^0(\bar\lambda).
\]
Therefore $\mu^0(\bar\lambda)=-1$, and uniqueness gives
$\bar\lambda=\lambda_*^0$.  Since every convergent subsequence has
the same limit, the whole family converges:
\[
\lambda_*^\varepsilon\longrightarrow\lambda_*^0.
\]
This proves \eqref{eq:root-convergence} and completes the proof.
\end{proof}

\end{document}
